%----------------------------------------------------------------------------------------
%	                            PORTADA
%----------------------------------------------------------------------------------------
% \makeportada{}{Daniel Carbajales Soto}{ 1° CFGS Automatización y Robótica Industrial}{}

%----------------------------------------------------------------------------------------
%                                    HEADER 
%----------------------------------------------------------------------------------------
% \header{}{Daniel Carbajales Soto}{1° CFGSARI}
%----------------------------------------------------------------------------------------
%                               TABLE OF CONTENTS
%----------------------------------------------------------------------------------------
% \maketoc
%----------------------------------------------------------------------------------------
%	                            TABLES
%----------------------------------------------------------------------------------------
% this is just a demo in case i need to make a table
%
%
%
% \begin{tabularx}{\linewidth}{|C{2cm}|Y|Y|C{2cm}|}
% \hline
% \makecell{Header 1} & \makecell{Header 2 \\ with linebreak} & \makecell{Header 3} & \makecell{Header 4} \\
% \hline
% Row 1 & This is a long cell text that automatically wraps in the Y column & Another long cell & Short \\
% \hline
% Row 2 & \multirow{2}{*}{Multirow text} & More text & Short \\
% \cline{1-1} \cline{3-4}
% Row 3 & & Extra text & Short \\
% \hline
% \end{tabularx}


%----------------------------------------------------------------------------------------
%	                            DOCUMENT
%----------------------------------------------------------------------------------------
\clearpage
\fchapter{5}{Ej 1 Pág. 27.}
\fsection{\boldit{INVESTIGACIÓN Y EXPRESIÓN ORAL.}}Realizad una investigación en grupo sobre las diferencias entre digitalización y transformación digital, utilizando ejemplos específicos para demostrar cómo se aplican en los siguientes sectores industriales: tecnología e informática, salud y farmacéutica, finanzas y banca, energía y recursos naturales, manufactura y producción industrial, comercio electrónico, telecomunicaciones, construcción, educación y formación, transporte y logística.

\fsubsection{Diferencias entre Digitalización y Transformación Digital}
La \textbf{digitalización} se refiere al proceso de convertir información o procesos analógicos en formato digital, con el objetivo principal de mejorar la eficiencia operativa y reducir costos en tareas específicas. Por ejemplo, implica escanear documentos en papel para almacenarlos en un sistema electrónico, lo que facilita el acceso y la gestión, pero no altera fundamentalmente la estructura de la organización.En contraste, la \textbf{transformación digital} es un enfoque estratégico más amplio que integra tecnologías digitales (como IA, big data o cloud computing) en todos los aspectos de una empresa, cambiando su modelo de negocio, cultura organizacional y experiencia del cliente para generar valor innovador y adaptabilidad a largo plazo.Mientras la digitalización es un paso inicial y táctico, la transformación digital es holística y disruptiva, requiriendo un cambio cultural profundo.A continuación, se analizan estas diferencias con ejemplos específicos en cada sector industrial solicitado. Para cada uno, se incluye un caso de digitalización (enfoque en conversión y optimización) y uno de transformación digital (enfoque en innovación estratégica).

\fsubsection{Tecnología e Informática}
\begin{itemize}
\item \textbf{Digitalización}: Una empresa de software escanea manuales impresos de sus productos y los convierte en PDFs accesibles en línea, reduciendo el tiempo de búsqueda de información para los usuarios.\item \textbf{Transformación Digital}: Microsoft integra IA en sus herramientas como Azure para crear ecosistemas de nube colaborativos, permitiendo a las empresas desarrollar aplicaciones personalizadas que transforman sus operaciones en tiempo real, generando nuevos modelos de ingresos basados en suscripciones.
\end{itemize}

\fsubsection{Salud y Farmacéutica} 
\begin{itemize}
\item \textbf{Digitalización}: Una farmacéutica digitaliza registros de pacientes en papel, almacenándolos en bases de datos electrónicas para agilizar consultas médicas y reducir errores de transcripción.
\item \textbf{Transformación Digital}: Pfizer utiliza IA y blockchain en su cadena de suministro para rastrear medicamentos en tiempo real, optimizando la producción flexible y personalizando tratamientos basados en datos genómicos, lo que acelera el desarrollo de vacunas como la de COVID-19.
\end{itemize}

\fsubsection{Finanzas y Banca}
\begin{itemize}
\item \textbf{Digitalización}: Un banco escanea cheques y contratos en papel para procesarlos digitalmente, acelerando las transacciones y minimizando el almacenamiento físico.
\item \textbf{Transformación Digital}: BBVA implementa plataformas de banca abierta con API y machine learning para ofrecer servicios personalizados como préstamos predictivos basados en datos del cliente, transformando el modelo de negocio hacia fintech colaborativas y expandiendo a ecosistemas no tradicionales como pagos móviles.
\end{itemize}

\fsubsection{Energía y Recursos Naturales}
\begin{itemize}
\item \textbf{Digitalización}: Una compañía petrolera digitaliza mapas geológicos en papel para bases de datos GIS, facilitando el análisis inicial de reservas sin alterar procesos de extracción.
\item \textbf{Transformación Digital}: Enel utiliza IoT y big data en redes inteligentes para integrar energías renovables, prediciendo demandas y optimizando la distribución en tiempo real, lo que reduce emisiones y crea modelos de negocio basados en prosumer (productores-consumidores).
\end{itemize}

\fsubsection{Manufactura y Producción Industrial}
\begin{itemize}
\item \textbf{Digitalización}: Una fábrica convierte registros manuales de inventario en hojas de cálculo digitales, mejorando el seguimiento de stock sin cambiar la línea de producción.
\item \textbf{Transformación Digital}: Siemens aplica el concepto de Industria 4.0 con robótica colaborativa y gemelos digitales, permitiendo producción personalizada a escala y mantenimiento predictivo, lo que revoluciona la cadena de valor al integrar proveedores en una red global conectada.
\end{itemize}

\fsubsection{Comercio Electrónico}
\begin{itemize}
\item \textbf{Digitalización}: Una tienda online digitaliza catálogos impresos en formatos web estáticos, facilitando búsquedas básicas de productos.
\item \textbf{Transformación Digital}: Amazon emplea IA en recomendaciones personalizadas y logística predictiva, transformando el e-commerce en un ecosistema de suscripciones (Prime) y fulfillment centers automatizados, que integra ventas físicas y digitales para una experiencia omnicanal.
\end{itemize}

\fsubsection{Telecomunicaciones}
\begin{itemize}
\item \textbf{Digitalización}: Una operadora telecom digitaliza facturas en papel para envíos por email, reduciendo costos de impresión y distribución.
\item \textbf{Transformación Digital}: Telefónica integra 5G y edge computing para ofrecer servicios como realidad aumentada en redes, cambiando el modelo de negocio hacia plataformas de datos que habilitan IoT masivo y colaboraciones con industrias como el gaming o la salud remota.
\end{itemize}

\fsubsection{Construcción}
\begin{itemize}
\item \textbf{Digitalización}: Una constructora escanea planos en papel para archivos CAD digitales, agilizando revisiones sin modificar flujos de trabajo en sitio.
\item \textbf{Transformación Digital}: Skanska usa BIM (Building Information Modeling) y drones con IA para modelado 3D colaborativo en tiempo real, integrando datos de sostenibilidad y permitiendo construcciones modulares predictivas que reducen desperdicios y generan nuevos servicios de consultoría digital.
\end{itemize}

\fsubsection{Educación y Formación}
\begin{itemize}
\item \textbf{Digitalización}: Una universidad convierte libros de texto en PDFs descargables, facilitando el acceso remoto a materiales sin cambiar el formato de clases.
\item \textbf{Transformación Digital}: Coursera integra IA en plataformas de aprendizaje adaptativo, personalizando cursos basados en datos de progreso y colaborando con empresas para certificaciones laborales, transformando la educación en un modelo híbrido global que prioriza habilidades sobre diplomas tradicionales.
\end{itemize}

\fsubsection{Transporte y Logística}
\begin{itemize}
\item \textbf{Digitalización}: Una empresa de logística digitaliza manifiestos de carga en papel para sistemas de rastreo básicos, mejorando la visibilidad inmediata de envíos.
\item \textbf{Transformación Digital}: DHL utiliza blockchain y IA en su red de supply chain para optimizar rutas predictivas y automatizar aduanas, creando un ecosistema conectado que integra vehículos autónomos y reduce emisiones, evolucionando hacia servicios de logística como servicio (LaaS).
\end{itemize}
En resumen, la digitalización es un prerrequisito esencial para la transformación digital, pero esta última impulsa la innovación competitiva. Los sectores que avanzan en transformación digital no solo sobreviven, sino que lideran cambios disruptivos, como se evidencia en los ejemplos.

%%%%%%%%%%%%%%%%%%%%%%%%%%%%%%%%%%%%%%%%%%%%%%%%%%%%%%%%%%
%%%%%%%%%%%%%%%%%%%% bibliography %%%%%%%%%%%%%%%%%%%%%%%%
% \nocite{*}                                 %%%%%%%%%%%%%
% \printbibliography                         %%%%%%%%%%%%%
%%%%%%%%%%%%%%%%%%%%%%%%%%%%%%%%%%%%%%%%%%%%%%%%%%%%%%%%%%
%%%%%% Last page in case there is no bibliography %%%%%%%%
% \label{Last Page}                           %%%%%%%%%%%%%
%%%%%%%%%%%%%%%%%%%%%%%%%%%%%%%%%%%%%%%%%%%%%%%%%%%%%%%%%%

